% Part: incompleteness
% Chapter: theories-computability
% Section: first-incompleteness

\documentclass[../../../include/open-logic-section]{subfiles}

\begin{document}

\olfileid[fr]{inc}{tcp}{inc}

\olsection{$\Th{Q}$ n'admet aucune extension complète, cohérente et axiomatisable}

\begin{thm}
\ollabel{thm:first-incompleteness}
Il n'existe aucune extension de~$\Th{Q}$ qui soit
à la fois complète, cohérente et !!{axiomatizable}.
\end{thm}

\begin{proof}
Nous savons déjà qu'aucune extension cohérente
de~$\Th{Q}$ n'est décidable. Or toute théorie
complète et !!{axiomatizable} est décidable.
\end{proof}

\begin{explain}
Ce théorème est proche de la formulation originale,
due à Gödel en 1931, du premier théorème
d'incomplétude. Outre une terminologie plus moderne,
les différences essentielles sont les suivantes :
Gödel supposait la « $\omega$-cohérence » plutôt que
la simple « cohérence »; il ne pouvait pas non plus
parler de théorie « !!{axiomatizable} » en toute
généralité, car la notion formelle de calculabilité
n'existait pas encore. Les modèles formels du calcul
ont été développés pendant la décennie suivante,
notamment par Gödel, en grande partie pour
caractériser les théories auxquelles s'applique
le phénomène d'incomplétude.

Le théorème affirme que la complétude, la cohérence
et l'existence d'une axiomatisation calculable ne
peuvent être réunies. Chacune peut cependant être
abandonnée tout en conservant les deux autres :
$\Th{Q}$ est cohérente et axiomatisée par un ensemble
calculable, mais elle n'est pas complète; la théorie
incohérente est complète et admet une axiomatisation
calculable, par exemple $\{\eq/[0][0]\}$, mais elle
n'est pas cohérente; enfin, l'ensemble des
!!{sentence}s arithmétiques vrais est complet et
cohérent, mais il n'admet pas d'axiomatisation
calculable.
\end{explain}

\end{document}
